\documentclass{article}

\usepackage{amsmath,amsfonts,amssymb,amsthm,mathtools,thmtools,cleveref}
\usepackage{stmaryrd} % double bracket [[1, n]]
\usepackage{graphicx,dsfont,float,booktabs,multirow}

\theoremstyle{definition}
% \theoremstyle{plain}
\newtheorem{eg}{Example}
\newtheorem{define}{Definition}
\newtheorem{assump}{Assumption}
\newtheorem{proposition}{Proposition}
\newtheorem{lemma}{Lemma}
\newtheorem{theorem}{Theorem}
\newtheorem*{remark}{Remark}
\newtheorem{remarknum}{Remark}

%% some (math) macros
\DeclareMathOperator{\interior}{int}
\DeclareMathOperator*{\argmax}{arg\,max}
\DeclareMathOperator*{\argmin}{arg\,min}
\newcommand\halfopen[2]{\ensuremath{[#1,#2)}}
\newcommand\openhalf[2]{\ensuremath{(#1,#2]}}
\newcommand*\diff{\mathop{}\!\mathrm{d}}

\setlength{\skip\footins}{15pt}

\begin{document}

\section*{Reachability analysis}

Define the set of admissible controls
\begin{align}
  \mathcal{U} = \{ u : [t_0, t_1] \rightarrow U \; \mid \; u \text{ is measurable} \} .
\end{align}
Define the forward reachable set
\begin{align}
  \mathcal{R}^f(t) = \mathcal{R}^f(t ; t_0, x_0) = \{ x(t ; t_0, x_0, u) \mid u \in \mathcal{U} \} ,
\end{align}
where $x(\cdot \, ; t_0, x_0, u)$ denotes the unique trajectory satisfying
\begin{align}
  \dot{x}(t) = Ax(t) + Bu(t), \quad x \in K, \quad  x(t_0) = x_0 ,
\end{align}
where $x \in K$ represents the state-constraints.
%
Similarly, define the backward reachable set
\begin{align}
  \mathcal{R}^b(t) = \mathcal{R}^b(t ; t_1, x_1) = \{ x_0 \mid \exists \, u \in \mathcal{U} \text{ s.t. } x(t_1 ; t, x_0, u) = x_1 \} .
\end{align}
%
Define $\mathcal{R}(t) = \mathcal{R}^f(t; t_0, x_0) \cap \mathcal{R}^b(t; t_1, x_1)$.

\paragraph{Some general questions.}

\begin{itemize}
  \item $\mathcal{R}^f(t)$ is compact (Filippov, extended to state constraints)
  \item $\mathcal{R}^b(t)$ is compact (time reversal + Filippov, extended to state constraints)
  \item $\mathcal{R}(t)$ is convex, so connected
\end{itemize}
Fillipov's theorem as mentioned in Liberzon's book does not handle state constraints.
It seems that ``viability theory'' could provide an answer to this question.

Let $\pi : \mathbb{R}^2 \rightarrow \mathbb{R}$ be the projection that maps each
state $q = (x, v)$ to its position component $\pi(q) = x$.
\begin{itemize}
  \item $\pi(\mathcal{R}(t))$ is an interval
  \item $t \mapsto \sup \pi(\mathcal{R}(t))$ with $v(t) = \dot{p}(t)$ is an optimal
        trajectory
\end{itemize}



\paragraph{Reachable set of time-space.}
Let $d \in \{\mathrm{forward}, \mathrm{backward} \}$, then the $d$ reachable set
over an interval to be
\begin{align}
  \mathcal{R}^d([t_0, t_1]) := \bigcup_{t \in [t_0, t_1]} (t, \mathcal{R}^d(t)) .
\end{align}
%
If $\mathcal{R}([t_0, t_1]) := \mathcal{R}^f([t_0, t_1]; x_0) \cap \mathcal{R}^b([t_0, t_1]; x_1) = \varnothing$,
then there does not exist any admissible control that steers the system from
state $x_0$ at time $t_0$ to state $x_1$ at time $t_1$.

\newpage

\begin{define}
  Let $q_0 = (x_0, v_0)$ be some initial state and consider some fixed control
  input $\omega \in [\omega_d, \omega_a]$, then we define
  \begin{align}
    x_{q_0}^\omega(t) := x_0 + \int_0^t \{ v_0 + \omega \tau \}_{[0, 1]} \diff \tau ,
  \end{align}
  where ${\{ \alpha \}}_{[0,1]} := \max(0, \min(1, \alpha))$ denotes clipping to
  interval $[0, 1]$.
\end{define}

\begin{proposition}
  The reachable set of positions at time $t$ is the interval
  \begin{align}
    \pi(\mathcal{R}^f(t)) = [x^{\omega_d}_{q_0}(t), x^{\omega_a}_{q_0}(t)] .
  \end{align}
\end{proposition}
\begin{proof}
  $\omega \mapsto x_{q_0}^\omega(t)$ is continuous.
\end{proof}

\end{document}
